\documentclass[12pt,a4paper]{article}

\usepackage{amsmath,amssymb}
\usepackage{graphicx}
\usepackage{geometry}
\usepackage{hyperref}
\usepackage{float}
\usepackage{natbib}
\usepackage{kotex}

\geometry{margin=1in}

\title{005\_BIMOD: 4-지점 이산 극성 분류 방법론}
\author{
윤덕진 (Solvaram) \\
\textit{팔체질 연구 플랫폼} \\
\texttt{with methodological clarification by DeepSeek AI}
}
\date{Version 3.2 -- 2026년 2월 24일}

\begin{document}

\maketitle

\begin{abstract}
본 논문은 BIMOD 프레임워크의 이산적 방법론적 기초를 정식화한다. 
나시온 중심 기준 구조는 네 개의 해부학적 지점을 정의한다. 
각 지점은 이진 두개 극성 상태(+1 또는 -1)를 가지며, 8-상태 분류 체계를 형성한다. 
이 모델은 상위 수준의 물리적 가설과 독립적인 이산적 구조적 분류 체계로 제시된다.
\end{abstract}

\section{방법론적 범위}

005\_BIMOD 논문은 분류 아키텍처만을 정의한다. 
연속적 벡터장, 자기 메커니즘, 생물물리학적 공명 모델을 주장하지 않는다. 
이러한 상위 수준 해석은 BIMOD-500에서 다루어지며, 여기서 확립된 분류 기반 위에 물리적 가설을 통합한다.

\section{구조적 정의}

나시온을 해부학적 기준 원점으로 정의한다. 이는 개인 간 위치적 안정성으로 알려진 두개계측 랜드마크이다 \citep{nasion_wikipedia}.

나시온 주위에 대칭적으로 네 개의 지점을 정의한다:

\[
L = \{L_1, L_2, L_3, L_4\}
\]

각 지점은 이진 극성 상태를 갖는다:

\[
P(L_j) \in \{+1,-1\}
\]

분류 공간은 데카르트 곱으로 정의된다:

\[
\mathcal{C} = L \times \{+1,-1\}
\]

따라서 총 분류 상태의 수:

\[
|\mathcal{C}| = 4 \times 2 = 8
\]

이 이산적 구조는 범주적 배타성을 보장한다: 각 개인은 여덟 가지 가능한 상태 중 정확히 하나에 해당한다.

\section{기하학적 참조}

\begin{figure}[H]
\centering
\includegraphics[width=0.7\textwidth]{005_NasionAI.jpg}
\caption{
네 개의 측정 지점의 대략적 위치를 보여주는 전체 얼굴 나시온 참조 이미지. 
각 지점(L1–L4)은 두 개의 극성 상태를 가지며, 8-상태 분류 체계를 형성한다. 
실제 검출 절차는 시각적 검사가 아닌 BIMOD 프로토콜을 따른다.
}
\label{fig:nasion}
\end{figure}

\section{레이블 층과 선택적 매핑}

여덟 가지 상태는 선택적으로 체질 레이블에 매핑될 수 있다:

\begin{itemize}
    \item 1S (목양체질), 1N (금양체질)
    \item 2N (목음체질), 2S (금음체질)
    \item 3N (토음체질), 3S (수음체질)
    \item 4N (토양체질), 4S (수양체질)
\end{itemize}

그러나 이러한 해석적 매핑은 상위 수준의 이론적 층(BIMOD-500의 Track B)에 속하며, 이 방법론적 정의에 내재적이지 않다. 분류 구조는 특정 해석과 독립적으로 존재한다.

\section{통계적 재현성}

분류 재현성은 이항 프레임워크를 사용하여 평가할 수 있다 \citep{binomialtest}:

\[
P(X \ge k | N, p=0.5) = \sum_{i=k}^{N} \binom{N}{i} (0.5)^N
\]

\(N=20\)회 시행에 대해 유의수준 \(\alpha=0.05\)(단측)에서의 임계값은 \(k=15\)이며, \(P(X \ge 15) = 0.0207\)이다.

경험적 검증은 다음을 요구한다:
\begin{itemize}
    \item 맹검 프로토콜 (시행자가 샘플 극성을 모름)
    \item RAW 이미지 보존 (AI 보정이나 필터링 없음)
    \item 여러 시행자에 걸친 독립적 재현
\end{itemize}

\section{물리적 가설과의 관계}

생체자성나노입자는 생물학적 시스템에서 보고되었으며 \citep{Kirschvink1992}, 지자기 상호작용에 기반한 자기수용 모델이 제안되었다 \citep{Ritz2000}. 그러나 여기서 정의된 분류 체계는 특정 물리적 메커니즘에 의존하지 않는다. 이는 다양한 설명 모델을 수용할 수 있는 구조적 프레임워크 역할을 한다.

\section{경계 조건}

본 논문은 이산적 분류 체계만을 정의한다. 
물리적 설명 모델, 연속적 장 형식화, 경험적 응용은 BIMOD-500 및 시리즈의 다른 논문에서 다루어진다.

\section{결론}

나시온 기반 여덟 극성 체계는 이론적 해석과 경험적 검토에 열려 있으면서도 그로부터 독립적인 이산적 방법론적 분류 프레임워크로 제시된다.

\begin{thebibliography}{99}

\bibitem{Kirschvink1992}
Kirschvink, J. L. et al. (1992). Magnetite in human tissues: A mechanism for biomagnetic detection? \textit{Bioelectromagnetics}.

\bibitem{Ritz2000}
Ritz, T. et al. (2000). A model for photoreceptor-based magnetoreception in birds. \textit{Biophysical Journal}, 78(2), 707–718.

\bibitem{binomialtest}
Clopper, C. J. \& Pearson, E. S. (1934). The use of confidence or fiducial limits illustrated in the case of the binomial. \textit{Biometrika}, 26, 404–413.

\bibitem{nasion_wikipedia}
Nasion. Wikipedia. Accessed 2026.

\end{thebibliography}

\end{document}